\subsection{El hogar dentro del análisis multidimensional de seguridad}
Infosegura plantea que el control del crimen debe acompañarse de oportunidades
de desarrollo y participación ciudadana. Su agenda de prevención recomienda seguir
avisos de violencia intrafamiliar para prevenir feminicidios y actuar en los
espacios públicos de mayor uso. Esta perspectiva permite incorporar a MICS
como una fuente sobre relaciones y prácticas familiares, complementaria a
la medición de seguridad ciudadana \citep{infosegura2021}.

El vínculo analítico se establece por mecanismos que requieren evaluación:
formas de crianza, aceptación normativa del uso de fuerza, autonomía de las
adolescentes y condiciones materiales del hogar. Son hipótesis para orientar
investigación y prevención. Los agregados disponibles no permiten saber si
los niños disciplinados violentamente son quienes sufrirán o ejercerán
violencia posterior, ni si las mujeres que justifican un escenario son las
víctimas de un feminicidio registrado años después.

La articulación exige observar también el entorno. La posibilidad de usar
transporte o espacios comunitarios y el acceso al empleo pueden relacionarse
con autonomía y acceso a servicios; no se estima aquí la dirección ni magnitud
de esas relaciones. En particular, una unión iniciada antes de los 18 años
plantea preguntas sobre agencia y oportunidades, pero no identifica por sí
sola la ocurrencia de cada tipo de violencia de pareja.

\begin{table}[htbp]
\centering\small
\caption{Correspondencia entre MICS e Infosegura}\label{tab:comparabilidad}
\begin{tabularx}{\textwidth}{@{}p{2.5cm}YYp{3.4cm}@{}}\toprule
Fuente y dimensión & Población y período & Unidad & Relación con el estudio\\\midrule
MICS: disciplina & Niño seleccionado 1-14, último mes en 2014 & Porcentaje con actos declarados & Crianza y protección infantil\\
MICS: normas y unión & Mujeres 15-49; cohorte 20-24; 2014 & Opiniones y eventos previos de unión & Género y autonomía\\
Infosegura: homicidios & Población total o por sexo; 2013-2021 & Tasas por cien mil; 2021 proyectado & Contexto de violencia letal\\
Cultura de Paz & Población encuestada; 2017-2021 & Victimización en 12 meses y percepción & Experiencia delictiva y espacios públicos\\
Registros y privaciones & Cortes propios de cada gráfico & Denuncias mensuales; variaciones en pp & Demanda institucional y contexto material\\\bottomrule
\end{tabularx}
\fuente{Cuestionario MICS y PNUD/Infosegura (2021). La serie de victimización no explicita la edad del universo; no se le atribuye la base 18+ de los indicadores de uso de espacios públicos.}
\end{table}
